\section{Specification: Wire Format, Gas, ABI}
\label{sec:specification}

\subsection{Slot and config key}

\begin{table}[h]
\centering
\small
\begin{tabular}{@{}ll@{}}
\toprule
\textbf{Attribute} & \textbf{Value} \\
\midrule
EVM precompile address & \texttt{0x0000000000000000000000000000000000012220} \\
Go address symbol & \texttt{ContractStarkFRIVerifyAddress} \\
Chain-config key & \texttt{starkfriVerify} \\
Go package & \texttt{starkfri} \\
Reference path & \texttt{lux/precompile/starkfri/} \\
\bottomrule
\end{tabular}
\caption{Slot and registration identifiers. The slot pin in
LP-300~\cite{lp300} is normative; this table is informative.}
\label{tab:slot-id}
\end{table}

The chain-config key \texttt{starkfriVerify} is the unique key under
which the upgrade-config block is keyed in the node's chain
configuration JSON. The reference implementation registers under
this key in \texttt{module.go} via the standard precompile
registration init function.

\subsection{Wire format}

Single-call interface; no Solidity function selector. The wire bytes
are:

\begin{lstlisting}[language={}]
[ 1 byte  ] version            -- 0x01 (VersionV1)
[ 4 bytes ] proof_len           -- big-endian uint32
[ proof_len bytes ] proof       -- must begin with MagicHeader "P3Q1"
[ 4 bytes ] pub_len             -- big-endian uint32
[ pub_len bytes ] public_inputs
\end{lstlisting}

\begin{table}[h]
\centering
\small
\begin{tabular}{@{}llll@{}}
\toprule
\textbf{Field} & \textbf{Size} & \textbf{Encoding} & \textbf{Constraint} \\
\midrule
\texttt{version}        & 1 byte  & raw byte       & MUST equal \texttt{0x01} \\
\texttt{proof\_len}     & 4 bytes & big-endian \texttt{uint32} & MUST equal length of \texttt{proof} \\
\texttt{proof}          & variable & opaque bytes   & First 4 bytes MUST equal \texttt{"P3Q1"} \\
\texttt{pub\_len}       & 4 bytes & big-endian \texttt{uint32} & MUST equal length of \texttt{public\_inputs} \\
\texttt{public\_inputs} & variable & opaque bytes   & format dictated by circuit \\
\bottomrule
\end{tabular}
\caption{Wire-format field table. Big-endian for length fields
matches Ethereum-mainnet conventions and the LP-218 P3Q ABI.}
\label{tab:wire-fields}
\end{table}

\subsubsection{Minimum input length}

\begin{equation*}
\texttt{MinInputLength} = 1 + 4 + 4 + 4 = 13 \text{ bytes}
\end{equation*}

The thirteen bytes decompose as: 1 version byte $+$ 4 proof-length
bytes $+$ 4 magic-header bytes (which live inside the
\texttt{proof} blob but are structurally required) $+$ 4
public-input-length bytes. Calls shorter than \SI{13}{\byte} are
rejected with \texttt{ErrInvalidInputLength}.

The Jasmin reference verifier checks
\texttt{len(input) >= MinInputLength} as the first gate, before the
version byte or magic-header check. This ordering is intentional:
length checks are cheaper than byte-string equality on the magic
header, and the dudect harness measures the dispatch path under all
three rejection modes (short, wrong-version, wrong-magic) to
empirically confirm timing uniformity.

\subsubsection{Magic header}

The 4-byte tag \texttt{MagicHeader = "P3Q1"} is preserved verbatim
from the pre-rename era so EasyCrypt~/ Lean~/ Jasmin~/ dudect proof
artifacts remain byte-identical against on-chain calldata. The
header is a \emph{wire-bytes tag, not a name claim}; renaming would
invalidate the proof artifacts. A future v2 wire format may rename
it to \texttt{"SFR1"} once the proof artifacts are refreshed and
the artifact-budget tracker (\texttt{scripts/checks/ec-admits.sh})
is updated in lockstep. This is documented at
\texttt{contract.go:46-52} in the reference implementation.

\subsubsection{Version byte}

Only \texttt{VersionV1 = 0x01} is accepted today. Versions
\texttt{0x00} or $\ge \texttt{0x02}$ are rejected with
\texttt{ErrInvalidVersion}. The version byte exists so a future LP
can graft a v2 wire format (post-Plonky3-upstream-rotation or a
recursive variant) without re-allocating the slot.

\subsection{Gas schedule}

\begin{equation*}
\texttt{RequiredGas}(\texttt{input}) = \texttt{BaseVerifyGas} + |\texttt{input}| \cdot \texttt{PerByteGas}
                                       = 200{,}000 + |\texttt{input}| \cdot 10
\end{equation*}

\begin{table}[h]
\centering
\small
\begin{tabular}{@{}lrr@{}}
\toprule
\textbf{Constant} & \textbf{Value} & \textbf{Units} \\
\midrule
\texttt{BaseVerifyGas} & \num{200000} & gas \\
\texttt{PerByteGas}    & \num{10}     & gas / byte \\
\bottomrule
\end{tabular}
\caption{Gas constants per \texttt{contract.go:107-110}.}
\label{tab:gas-constants}
\end{table}

The \num{200000} base reflects FRI verify cost being
\emph{logarithmic in circuit size}: on-chain we only see the
serialized proof, so we charge a flat base plus per-byte for
bandwidth and decode. At a typical \SI{30}{\kilo\byte} proof, total
gas is $200{,}000 + 30{,}000 \cdot 10 = 500{,}000$ gas per verify --
$\sim$$24$~verifies per \num{12000000}-gas C-Chain block.

\subsubsection{Upfront gas billing}

Gas is charged upfront before any structural validation or backend
dispatch (\texttt{contract.go:194-198}). The
\texttt{oog\_dominates} EasyCrypt lemma
(\texttt{proofs/easycrypt/P3Q\_Gas\_Model.ec}) pins this:
out-of-gas wins against every other failure mode. The lemma's
statement is informally ``for every input and every
\texttt{gas\_in} value, if
\texttt{gas\_in < RequiredGas(input)} then the precompile's return
tuple is $(\texttt{nil}, 0, \texttt{ErrOutOfGas})$ regardless of any
other field's contents.''

This matters for soundness because it forecloses a gas-side channel:
a caller cannot probe which validation stage rejected the input by
arranging \texttt{gas\_in} to be just-below the per-stage threshold,
because there are no per-stage thresholds. There is one threshold,
charged upfront, and either it is met (in which case validation
runs) or it is not (in which case the call short-circuits with
\texttt{ErrOutOfGas} before any byte of the input is parsed).

\subsection{Return contract}

\begin{table}[h]
\centering
\small
\begin{tabular}{@{}p{6.5cm}p{6.5cm}@{}}
\toprule
\textbf{Outcome} & \textbf{Return} \\
\midrule
Verified & $(\texttt{nil}, \texttt{gasLeft}, \texttt{nil})$ \\
Insufficient gas & $(\texttt{nil}, 0, \texttt{contract.ErrOutOfGas})$ \\
Input shorter than \texttt{MinInputLength}, or internal length fields exceed remaining bytes & $(\texttt{nil}, \texttt{gasLeft}, \texttt{ErrInvalidInputLength})$ \\
\texttt{version != 0x01} & $(\texttt{nil}, \texttt{gasLeft}, \texttt{ErrInvalidVersion})$ \\
\texttt{proof[0:4] != "P3Q1"} & $(\texttt{nil}, \texttt{gasLeft}, \texttt{ErrInvalidProof})$ \\
No verifier registered via \texttt{RegisterVerifier} & $(\texttt{nil}, \texttt{gasLeft}, \texttt{ErrVerifierNotRegistered})$ \\
Backend returned \texttt{(false, nil)} & $(\texttt{nil}, \texttt{gasLeft}, \texttt{ErrInvalidProof})$ \\
Backend returned \texttt{(\_, err)} for \texttt{err != nil} & $(\texttt{nil}, \texttt{gasLeft}, \texttt{err})$ -- backend internal failure surfaces verbatim \\
\bottomrule
\end{tabular}
\caption{Return-tuple contract by outcome. Per
\texttt{contract.go:113-118}, error sentinels are exported package
symbols so Solidity callers (via revert-data decoding) and Go
callers can pattern-match on them.}
\label{tab:return-contract}
\end{table}

This contract differs from the LP-218 \texttt{abiFalse}-on-failure
convention because LP-221 reports \emph{structural} failures via
Go error returns. A Solidity caller will see the precompile call
revert on any structural mismatch (length, version, magic) and on
backend verification failure; success returns empty output (the
proof either verifies or it doesn't, and the caller's contract
treats the absence of revert as success).

The \texttt{ErrInvalidProof} sentinel is overloaded to cover both
\emph{wire-format} bad-magic and \emph{cryptographic} backend
rejection. A future LP may split these into two sentinels
(\texttt{ErrBadMagic} vs \texttt{ErrBackendReject}); today they
share the same revert-data path because the caller's response is
identical (reject the batch). Splitting them is a UX improvement,
not a soundness improvement; it lands behind the version-byte gate
when the wire format bumps.

\subsection{Cross-reference to the registry}

LP-300~\S``Post-quantum precompile registry'' is the master pin for
slot \texttt{0x012220}. The pre-LP-221 row in LP-300 read:

\begin{quote}
\texttt{0x012220 -- LP-218 -- Placeholder for displaced
Plonky3-fork STARK -- pending dedicated LP allocation.}
\end{quote}

LP-221 closes the placeholder. LP-300's row should be amended on
its next revision to:

\begin{quote}
\texttt{0x012220 -- LP-221 -- STARK-FRI -- Plonky3-fork FRI~/
cSHAKE256~/ Goldilocks STARK verifier; magic header "P3Q1"
preserved verbatim for proof-artifact stability -- Draft --
2026-06-04.}
\end{quote}

LP-300 is normative; LP-221 is the dedicated LP that LP-300's
placeholder row was waiting on.
